% CHAPTER 2 — PRODUCT SPECIFICATIONS

\chapter{Product Specifications}

\section{Introduction}
This chapter defines the full specification of the BodyQ project. It
covers the project management approach adopted by the team, the complete
product backlog, the sprint planning and execution, and the detailed
functional and non-functional requirements of the system. It concludes
with an analysis of the project's constraints, limitations, assumptions,
standards, and ethical considerations.

\section{Project Management}

\subsection{Methodology}
The BodyQ project was developed using an \textbf{Agile Scrum}
methodology. Scrum was chosen because the project required the integration
of several technically heterogeneous components---a mobile application,
an AI engine, a backend, and an administration dashboard---whose
interfaces and requirements were expected to evolve as implementation
progressed. Scrum's iterative sprint structure allowed the team to
deliver working increments at the end of each sprint, validate them
against expected behavior, and adjust the plan for the following sprint
based on what was learned.

For the Team of 5 students, sprint planning, review, and retrospective meetings were conducted at the beginning and end of each sprint. A
shared JIRA task board tracked user story status across To~Do,
In~Progress, and Done columns throughout the Bocklog while tracking progress and contributions. Sprint duration was two weeks, with the
exception of Sprint~0, which was one week long.

\subsection{Timeline}
The project was planned across five sprints over nine weeks, as detailed
in Table~\ref{tab:sprints}.

\begin{table}[H]
	\centering
	\caption{Sprints Planning}
	\label{tab:sprints}
	\begin{tabularx}{\textwidth}{lXcc}
		\toprule
		\textbf{Sprint} & \textbf{Primary Deliverables} & \textbf{From} & \textbf{To} \\
		\midrule
		Sprint~0 --- Foundation
		& Repository setup, Supabase project
		& Week~1\\[3pt]
		Sprint~1 --- Auth \& Profile
		& full DB schema, seed data, Authentication, 7-step onboarding, BMR/TDEE calculation, WebPage Initiation, Admin Dashboard implementation
		& Week~2 & Week~3 \\[3pt]
		Sprint~2 --- Core Tracking
		& Workout module, exercise library, nutrition logging, barcode scanner, water tracker, home dashboard
		& Week~3 & Week~4 \\[3pt]
		Sprint~3 --- AI Engine
		& Posture detection MVP, food recognition, RAG coaching engine
		& Week~5 & Week~6 \\[3pt]
		Sprint~4 --- Advanced Features
		& Weekly reports, correlation engine, habit/streak system, push notifications
		& Week~7 & Week~8 \\[3pt]
		Sprint~5 --- Polish \& Ship
		& Testing, bug fixes,UI animations, performance optimization, experimenting with  wearable device
		& Week~9 & Week~9 \\
		\bottomrule
	\end{tabularx}
\end{table}

\section{Product Backlog}
% Replace your entire table block with this:
\begin{xltabular}{\textwidth}{Xcc}
	\caption{Product Backlog} \label{tab:backlog} \\
	\toprule
	\textbf{User Story} & \textbf{Priority} & \textbf{Sprint} \\
	\midrule
	\endfirsthead
	
	% Header for Page 2+ (No "Continued" text)
	\toprule
	\textbf{User Story} & \textbf{Priority} & \textbf{Sprint} \\
	\midrule
	\endhead
	
	% Footer for all pages except the last (Empty)
	\endfoot
	
	% Footer for the very last page
	\bottomrule
	\endlastfoot
	
	% YOUR DATA START
	As a user, I want to create an account so that my data is saved securely. & High & S1 \\
	As a user, I want an easy to use dashboard, so I can see my daily progress overview. & High & S1 \\
	As a user, I receive personalized calorie and macro targets & High & S1 \\
	As a user, I want to define my age, height, weight, goals, and experience level so the system can personalise recommendations. & High & S1 \\
	As a user, I can browse and search the exercise library & High & S2 \\
	As a user, I want to receive reminders for workouts, hydration, and sleep, so i can be on track & High & S2 \\
	As a user, I can log my meals and track & High & S2 \\
	As a user, I want to track my daily steps, so that I can monitor my activity level. & High & S2 \\
	As a user, I want the app to record my workout duration and exercises performed, so that I can get a clear insight on my progress. & High & S3 \\
	As a user, I want to view a list of supported exercises with instructions, so that I can choose what suits me better & High & S3 \\
	As a user, I want to receive daily motivational quotes, so that I feel inspired to stay active. & High & S3 \\
	As a user, I want to be able to customise my dashboard view, so that it fits my preferences. & High & S3 \\
	As a user, I want to earn rewards for completing my fitness goals, so that I can stay motivated. & Medium & S2 \\
	As a user, I want to receive AI-generated health advice based on my habits. & Medium & S3 \\
	As a user, I want to log my daily water intake, so that I can get a clear insight on my progress. & Medium & S3 \\
	As a user, I want to log my sleep duration, so that I can get a clear insight on my progress. & Medium & S4 \\
	As a user, I want to log meals manually, so that I can track my calorie intake, so that I can get a clear insight on my progress. & Medium & S4 \\
	As a user, I want to scan food using my camera to estimate calories. & Medium & S4 \\
	As a user, I want to see calorie summaries and suggestions, so that I can fit my diet to my workout. & Medium & S4 \\
	As a user, I want to see my progress over time, so that I can get a clear insight on my progress. & High & S4 \\
	As a user, I want to track progress toward fitness goals, so that I can get a clear insight on my progress. & Medium & S4 \\
	As an admin, I can review AI recommendation logs & Low & S4 \\
	As a user, I want to enable my camera during workouts so the AI can analyse my posture & Low & S4 \\
	As a user, I want to receive feedback when my posture is incorrect. & Low & S4 \\
	% YOUR DATA END
	
\end{xltabular}


\section{Sprint Backlog}
\subsection{Sprint 1}
% TODO: Add Sprint 1 backlog figure

\subsection{Sprint 2}
% TODO: Add Sprint 2 backlog figure

\section{Requirements Specification}

\subsection{Functional Requirements}
% TODO: Use case scenarios + diagrams

\subsection{Non-Functional Requirements}
% TODO: Performance, scalability, usability metrics

\section{Project Considerations}

\subsection{Project Constraints}
\subsection{Project Limitations}
\subsection{Project Delimitations}
\subsection{Project Assumptions}
\subsection{Project Standards}
\subsection{Business, Social and Ethical Considerations}

\section{Conclusion}
